\paragraph{Del origen incidencial a la modificación reticular.}
\label{inc-ampl-reticulo}
La bandera marcada determina una posición mediante una operación
conjuntista equivariante:
\[
 o_*=(H_e\cap H_\varphi)
       \setminus(H^-_\alpha\cup H_{\rm APP})=\{1\}.
\]
Su función consiste en distinguir una componente entre las doce que
intervienen en la realización reticular. El código ternario completo
actúa como código de pegado en el módulo discriminante
$(A_2^*/A_2)^{12}$. Su preimagen define $N(C_W)$, y la
autodualidad, los pesos y la distancia del código se traducen,
respectivamente, en propiedades de integralidad y determinante, paridad
y cotas de norma de las clases de pegado. Aquí la información simbólica
del código cumple una función que el soporte aislado ya no desempeña.

La métrica $A_2$ y la cámara radial positiva se fijan en el desarrollo
reticular. Dentro de esa cámara, la congruencia necesaria para un vecino
par de índice tres tiene doce realizaciones de norma mínima $54$,
permutadas por el cambio de componente distinguida. El origen $o_*$
selecciona
\[
 v_\alpha=4\rho_1+\rho_2+\cdots+\rho_{12},\qquad
 v_\alpha^2=2(4^2+11)=54.
\]
El coeficiente $4$ pertenece a esa solución de la congruencia en la
cámara declarada. El subíndice $\alpha$ registra la procedencia
incidencial de la selección: su contenido es el marco que intervino en
la clausura y que ahora individualiza una dirección reticular.

El emparejamiento con $v_\alpha$ determina el subretículo
$N_{\alpha,0}$ de índice tres, y la adición de la clase
$v_\alpha/3$ produce el vecino $N_\alpha$ de
\eqref{exc:vecino}. Esta modificación conserva el rango y recupera
unimodularidad y paridad; la selección por congruencia elimina las raíces
anteriores, y las cotas locales excluyen raíces en las nuevas clases.
La identificación con Leech aparece después de esas verificaciones. La
realización binaria de las cinco regiones y esta construcción reticular
son dos prolongaciones del marco común: la segunda recibe directamente
el código ternario como dato de pegado.

